\scanpage{19}
\mspage{17}
tallogène ne donnerait pas les résultats qu'il faut, on pourrait ess\inserted{a}yer (au lieu de prendre des foncteurs dérivés dans la catégorie de \uline{tous} les Modules sur le site annelé cristallog\deleted{\unsure{raphique}}\inserted{ène}) de prendre des foncteurs dérivés dans la catégorie des seuls Modules quasi-cohérents sur le site cristallogène, i.e. dans la catégorie des cristaux de modules. \unsure{?}

\uline{Chapitre} 2 \uline{La cohomologie de De Rham est un cristal}.

2.1. L'affirmation du titre n'est pour l'instant qu'une hypothèse ou un voeux pieux, mais je suis convaincu qu'elle est essentiellement correcte. Comme je l'ai dit dans 1.17., il y a pour le moment deux indications dans cette direction:

a) Le travail de Washnitzer-Monsky. Deux grosses imperfections dans leur tapis: $1^{\mathrm{o}}$ ils n'obtiennent que des invariants cohomologiques \uline{locaux} sur leur variété lisse en caractéristique $p>0$, via leurs relèvements; pour avoir un invariant global, ils doivent se limiter à la dimension 1. On peut penser que cela tient à leur manque de familiarité avec les machines cohomologiques. $2^{\mathrm{o}}$ leurs invariants sont des (faisceaux de) vectoriels sur le corps des fractions de $W$, et non sur $W$, i.e. ils n'obtiennent que des invariants ``modulo isogénie''. Il semble bien, en effet, que leur démonstration d'invariance fait intervenir de façon \deleted{ass} essentielle des dénominateurs. Il n'est pas exclu, d'après cette indication, que dans le titre du Chapitre il faille remplacer ``cristal'' par ``isocristal''\inserted{(et ``est'' par ``définit'')}. Ce serait bien dommage, mais n'exclurait pas pour autant l'existence d'une bonne théorie cohomologique pour les cristaux, qui pour un morphisme propre et lisse et le ``cristal unité'' coïnciderait \uline{modulo isogénie} avec ce que donne De Rham. Le test décisif reste celui indiqué dans 1.17, savoir: donne-t-il un résultat positif seulement en caractéristique zéro, ou en toute caractéristique ?

b) L'existence des connexions de Gauss-Manin. \deleted{Comme} J'ai vérifié pour tout morphisme lisse $f:\deleted{\unsure{S'}}\inserted{X}\longrightarrow S$ l'existence d'une connexion canonique \inserted{(absolue)} sur les $\underline{R}^if_*$ au sens de De Rham relatif, ou plus correctement, sur le complexe
\margin{18}
